\chapter{Planteamiento del Problema}
\label{chap:planteamiento}

La modelación de series temporales en economías emergentes como la de
Honduras presenta un desafío dual: por un lado, la necesidad de
capturar con precisión las tendencias cuantitativas y, por otro, la
de interpretar los cambios abruptos de régimen provocados por la
volatilidad interna y los shocks externos.  Los modelos
econométricos tradicionales ---ARIMA, GARCH, VAR--- a menudo luchan
por reconciliar estas dos facetas, ofreciendo un ajuste preciso con
escasa interpretabilidad, o bien una descripción cualitativa sin el
rigor necesario para un pronóstico fiable \cite{Hamilton1994,
Box2015, Tsay2010, Lutkepohl2005}.

Este vacío en la literatura, particularmente palpable en el contexto
hondureño, exige enfoques híbridos que entrelacen el análisis
determinista con la inferencia estocástica.  Las series presentadas
en el Capítulo~\ref{chap:seleccionvariable} ---PIB y
combustibles--- constituyen casos paradigmáticos: su estructura de
regímenes múltiples y transiciones persistentes no se captura
adecuadamente con modelos lineales clásicos.

% ===================================================================
\section{Revisión de Literatura y Posicionamiento}
\label{sec:revision_literatura}
% ===================================================================

La presente sección sitúa la contribución de esta tesis en el
contexto de tres corrientes de la literatura: (i)~modelos de cambio
de régimen para series temporales, (ii)~métodos de estimación de
matrices de transición, y (iii)~computación de reservorio aplicada
a series económicas.

% --- Modelos de Cambio de Régimen ---
\subsection{Modelos de Cambio de Régimen}

El enfoque más influyente para modelar transiciones entre estados
económicos es el modelo \textit{Markov-Switching Autoregressive}
(MS-AR) de Hamilton~\cite{hamilton1989new, Hamilton1994}, en el cual
los parámetros de un proceso AR cambian según una cadena de Markov
con probabilidades de transición no observables.  La estimación se
realiza por máxima verosimilitud (MLE) mediante el algoritmo EM,
que requiere la evaluación de la verosimilitud completa por
filtrado recursivo.  A pesar de su elegancia teórica, el modelo
MS-AR presenta dos limitaciones prácticas relevantes para esta
tesis:
\begin{enumerate}
  \item La función de verosimilitud es multimodal, lo que hace que la
    convergencia del algoritmo EM dependa de la inicialización y no
    garantice un óptimo global~\cite{Hamilton1994}.
  \item La complejidad computacional crece exponencialmente con el
    número de estados $K$, limitando las aplicaciones con más de
    3--4 regímenes.
\end{enumerate}

Otros modelos de cambio de régimen incluyen:
\begin{itemize}
  \item \textbf{Threshold AR (TAR)} \cite{Tong1990, Tsay2005}: el
    régimen se determina por una variable umbral observable,
    eliminando la componente latente.  La función de transición es
    discontinua, similar a la función $\pi$ del modelo TCROC-Markov.
  \item \textbf{STAR/LSTAR} \cite{Terasvirta1994}: generalizan el
    TAR con una función de transición suave (logística o
    exponencial), permitiendo cambios graduales entre regímenes.
  \item \textbf{HMM gaussiano} \cite{rabiner1989tutorial}: los
    estados son latentes y las emisiones continuas.  Comparte la
    estructura estocástica con nuestro modelo, pero difiere en que
    el TCROC-Markov determina los estados \emph{explícitamente}
    a partir de la tasa de cambio relativa, no como variable
    oculta.
  \item \textbf{Modelos factoriales dinámicos}
    \cite{KimNelson1999}: combinan factores latentes con cambio de
    régimen.  Son más generales pero más difíciles de estimar.
\end{itemize}

La contribución del marco TCROC-Markov respecto a estos modelos es
la separación explícita entre la etapa de discretización
(determinista, basada en el operador TCROC) y la etapa de modelado
estocástico (estimación de la matriz de transición vía NNLS).  Esta
arquitectura de dos etapas evita la optimización no convexa del
MS-AR y produce estados con interpretación económica directa.

% --- Estimación de Matrices de Transición ---
\subsection{Estimación de Matrices de Transición}

Para cadenas de Markov con estados observados, el estimador de
máxima verosimilitud de la matriz de transición coincide con las
frecuencias relativas de transición
\cite{AndersonGoodman1957, Norris1998}.  Este estimador es
consistente y asintóticamente normal
(Proposición~\ref{prop:normalidad_asintotica}), pero no incorpora
información estructural ni restricciones de suavidad.

El enfoque NNLS adoptado en esta tesis
\cite{LawsonHanson1995, KimPark2011} reformula la estimación como
un problema de mínimos cuadrados con restricciones de
no-negatividad.  Sus ventajas frente al MLE son:
\begin{enumerate}
  \item \textbf{Convexidad global:} el problema NNLS es convexo,
    garantizando convergencia al óptimo global sin dependencia de
    la inicialización.
  \item \textbf{Regularización implícita:} la restricción de
    no-negatividad actúa como regularizador, produciendo
    estimaciones más estables con muestras pequeñas.
  \item \textbf{Extensibilidad:} la formulación matricial
    (Capítulo~\ref{cap:tcroc_hibrido_integrado}) permite
    incorporar restricciones adicionales de forma natural.
\end{enumerate}

% --- Computación de Reservorio ---
\subsection{Computación de Reservorio para Series Económicas}

Los modelos de computación de reservorio (\textit{Reservoir
Computing}, RC) fueron introducidos por Jaeger~\cite{jaeger2001echo}
con las Redes de Estado de Eco (\textit{Echo State Networks}, ESN)
y han demostrado capacidad universal de
aproximación~\cite{grigoryeva2018universal, GononOrtega2020}.  La
propiedad fundamental es el Estado de Eco (ESP): la dinámica interna
del reservorio ``olvida'' sus condiciones iniciales, garantizando
que la respuesta dependa solo de la entrada
\cite{yildiz2012, lukosevicius2009reservoir}.

La extensión SSRC (\textit{Stochastically Structured Reservoir
Computers}) propuesta por Banegas y
Vides~\cite{BanegasVides2025, banegas2025} integra la estructura de
Markov directamente en la arquitectura del reservorio, utilizando
matrices de transición estocásticas como pesos de conectividad.
Esta integración es la base teórica del
Capítulo~\ref{cap:sec_ssrc}, donde se demuestra que el modelo
TCROC-Markov es un caso particular de los SSRC (Teorema de
Inclusión).

% --- Gap en la literatura ---
\subsection{Identificación del Gap}

La revisión anterior revela un vacío en la intersección de tres
áreas: (i)~los modelos de cambio de régimen clásicos no aprovechan
tasas de cambio relativas observables para la discretización;
(ii)~la estimación NNLS no ha sido aplicada sistemáticamente a
matrices de transición económicas; (iii)~la conexión entre modelos
de Markov lineales y computación de reservorio no ha sido
formalizada.  Esta tesis llena simultáneamente estos tres gaps
mediante el marco TCROC $\to$ Markov $\to$ NNLS $\to$ SSRC.

La Tabla~\ref{tab:comparativa_modelos} resume las diferencias
clave entre el modelo propuesto y los principales competidores.

\begin{table}[htbp]
\centering
\caption{Comparación de modelos de cambio de régimen.}
\label{tab:comparativa_modelos}
\resizebox{\textwidth}{!}{%
\begin{tabular}{@{}lcccc@{}}
\toprule
\textbf{Criterio} &
\textbf{TCROC-Markov} &
\textbf{MS-AR} &
\textbf{TAR} &
\textbf{HMM} \\
\midrule
Estados &
Observados ($\pi$) &
Latentes (EM) &
Observados (umbral) &
Latentes (Baum-Welch) \\
Estimación de $P$ &
NNLS (convexo) &
MLE (no convexo) &
N/A &
MLE (no convexo) \\
Convergencia &
Global &
Local &
Global &
Local \\
Validez estocástica &
Por construcción &
No garantizada &
N/A &
No garantizada \\
Interpretabilidad &
Alta (regímenes) &
Media &
Alta &
Baja \\
Complejidad &
$O(K^2 T)$ &
$O(K^2 T)$ &
$O(KT)$ &
$O(K^2 T)$ \\
Extensión no lineal &
SSRC (Cap~\ref{cap:sec_ssrc}) &
No estándar &
STAR &
No estándar \\
\bottomrule
\end{tabular}%
}
\end{table}
% ===================================================================
\section[Hoja de Ruta Metodológica]{Hoja de Ruta Metodológica:
Del Análisis Determinista a la Inferencia Estocástica}
\label{sec:hoja_ruta}
% ===================================================================

La investigación se desarrolla en tres fases metodológicas
conectadas:

\textbf{Fase~1: Construcción del operador de discretización}
(Capítulo~\ref{cap:tcroc}).  Se formaliza la familia de
transformaciones TCROC y sus variantes paramétricas (TCROCM, ETCROC,
ETCROCM), demostrando rigurosamente las propiedades de existencia,
unicidad, consistencia asintótica y estabilidad ante
perturbaciones.  El
objetivo no es solo proponer un nuevo estimador de tasas, sino
construir un ``motor de discretización'' robusto e interpretable,
capaz de transformar una serie temporal continua en una secuencia de
regímenes económicos con significado empírico.

\textbf{Fase~2: Modelo estocástico híbrido}
(Capítulos~\ref{cap:TCROC-Markov}
y~\ref{cap:tcroc_hibrido_integrado}).  Se integra el operador TCROC
con cadenas de Markov para modelar las transiciones probabilísticas
entre regímenes.  La novedad central es la aplicación de un estimador
de Mínimos Cuadrados No Negativos (NNLS) para la matriz de
transición, garantizando por construcción la validez estocástica de
los parámetros ---un desafío persistente en modelos de cambio de
régimen como Markov-Switching \cite{Hamilton1994}.

\textbf{Fase~3: Validación empírica y extensiones}
(Capítulos~\ref{cap:regimenes_combustibles}
y~\ref{cap:sec_ssrc}).  Se valida el marco completo sobre los
precios semanales de combustibles hondureños (2017--2025), incluyendo
optimización exhaustiva de hiperparámetros, análisis espectral y
comparación con benchmarks.  Además, se propone una extensión no
lineal mediante Computación de Reservorio Estructurado (SSRC) que
generaliza el marco lineal.

% ===================================================================
\section{Objetivos y Alcance de la Investigación}
\label{sec:objetivos}
% ===================================================================

El objetivo primordial de este trabajo es entregar una metodología
completa, desde su concepción teórica hasta su validación empírica,
para el análisis de series temporales económicas y financieras.  Los
objetivos específicos son:

\begin{enumerate}
    \item Desarrollar formalmente el operador TCROC y sus variantes,
    con demostraciones rigurosas de sus propiedades matemáticas
    (Capítulo~\ref{cap:tcroc}).

    \item Integrar el operador con cadenas de Markov y estimar la
    matriz de transición mediante NNLS, asegurando la validez
    estocástica (Capítulos~\ref{cap:TCROC-Markov}
    y~\ref{cap:tcroc_hibrido_integrado}).

    \item Validar empíricamente el modelo en series de combustibles
    hondureños, comparándolo con cuantiles fijos y MS-AR
    (Capítulo~\ref{cap:regimenes_combustibles}).

    \item Evaluar la extensión no lineal del marco mediante SSRC
    (Capítulo~\ref{cap:sec_ssrc}).

    \item Proporcionar implementaciones computacionales replicables
    para facilitar la adopción del método.
\end{enumerate}

El beneficio derivado es doble: por un lado, se aporta al
conocimiento en modelación matemática una herramienta innovadora; por
otro, se ofrece a los tomadores de decisiones en los ámbitos
gubernamental y empresarial una metodología más precisa para la
planificación estratégica en entornos volátiles.

% ===================================================================
\section{Estructura del Documento}
\label{sec:estructura}
% ===================================================================

La tesis se organiza como sigue.  El
Capítulo~\ref{chap:seleccionvariable} presenta las series de
estudio.  El Capítulo~\ref{cap:tcroc} desarrolla la teoría del
operador TCROC.  El Capítulo~\ref{cap:TCROC-Markov} establece la
integración con cadenas de Markov.  El
Capítulo~\ref{cap:tcroc_hibrido_integrado} construye el modelo
híbrido NNLS.  El Capítulo~\ref{cap:regimenes_combustibles} valida
empíricamente el marco en combustibles.  El
Capítulo~\ref{cap:sec_ssrc} extiende el modelo con computación de
reservorio.  Finalmente, el Capítulo~\ref{cap:conclusiones}
sintetiza las conclusiones y las líneas de trabajo futuro.

La Figura~\ref{fig:diagrama_tesis} resume esquemáticamente esta
estructura y las dependencias entre capítulos.

\begin{figure}[htbp]
\centering
\begin{tikzpicture}[
  node distance=1.0cm and 1.0cm,
  block/.style={rectangle, draw=nordGray, rounded corners=3pt,
    fill=nordSnow, text width=2.8cm, text centered,
    minimum height=0.9cm, font=\footnotesize, text=nordDark},
  arrow/.style={-{Stealth[length=2.5mm]}, thick, nordGray}
]
% Fila 1
\node[block] (cap4) {Cap.~4\\Operador TCROC};
\node[block, right=1.0cm of cap4] (cap5) {Cap.~5\\TCROC + Markov};
\node[block, right=1.0cm of cap5] (cap8) {Cap.~8\\Extensión SSRC};
% Fila 2
\node[block, below=of cap5] (cap6) {Cap.~6\\Modelo NNLS};
% Fila 3
\node[block, below=of cap6] (cap7) {Cap.~7\\Validación Empírica};
% Fila 4
\node[block, below=of cap7] (cap9) {Cap.~9\\Conclusiones};
% Flechas
\draw[arrow] (cap4) -- (cap5);
\draw[arrow] (cap5) -- (cap6);
\draw[arrow] (cap5) -- (cap8);
\draw[arrow] (cap6) -- (cap7);
\draw[arrow] (cap7) -- (cap9);
\draw[arrow] (cap8) |- (cap9);
\end{tikzpicture}
\caption{Estructura de la tesis y dependencias entre capítulos.}
\label{fig:diagrama_tesis}
\end{figure}

Todo el código fuente, los datos y los scripts de reproducción
están disponibles en el repositorio público del proyecto\footnote{%
\url{https://github.com/NORSAB/EcoSeries-RegimeSwitching}.
DOI: \href{https://doi.org/10.5281/zenodo.18752541}{10.5281/zenodo.18752541}.}.
